\documentclass[../../../include/open-logic-section]{subfiles}

\begin{document}

\olfileid{pt}{nat}{qua}

\olsection{正则推导与代入}

% Part: proof-theory
% Chapter: natural-deduction
% Section: quantifiers

量词规则由于适用于 \Elim{\lexists} 和~\Intro{\lforall} 规则的“本征变元条件”而带来了一些复杂性。
这些复杂性大多可以通过确保这些推理中的常项符号~$c$ 不被重复使用来处理。我们引入以下定义。

\begin{defn}
一个自然演绎推导~$\delta$ 是\emph{正则的}，如果
\begin{enumerate}
  \item 没有本征变元~$a$ 被用作多于一个 \Elim{\lexists} 或~\Intro{\lforall} 推理的本征变元，并且
  \item 如果 $a$ 是某个 \Elim{\lexists} 或~\Intro{\lforall} 推理的本征变元，那么它仅分别出现在导出该推理的次前提或前提的直接子推导中。
\end{enumerate}
\end{defn}

\begin{lem}\ollabel{lem:subst-regular} 如果 $\delta$ 是正则的且以~$!C$ 结束，$c$ 是一个在~$\delta$ 中不作为本征变元使用的常项符号，且 $t$~是一个不包含~$\delta$ 的任何本征变元的项，那么将~$\delta$ 中每次出现的~$c$ 替换为~$t$ 所得到的 $\Subst{\delta}{t}{c}$ 是一个正则推导。$\Subst{\delta}{t}{c}$ 的末端公式是 $\Subst{!C}{t}{c}$。如果 $!A^x$ 是~$\delta$ 的一个开放假设，那么 $\Subst{!A}{t}{c}^x$ 是 $\Subst{\delta}{t}{c}$ 的一个开放假设。
\end{lem}

\begin{proof}
  对 $\pheight{\delta}$ 施归纳。如果 $\pheight{\delta} = 0$，则它仅由一个假设~$!A^x$ 构成。此时 $\Subst{\delta}{t}{c}$ 就是~$\Subst{!A}{t}{c}^x$。

  如果 $\pheight{\delta} > 0$，我们根据~$\delta$ 的最后一个推理分情况讨论。命题联结词规则的情况较为简单，留作练习。我们考虑 $\delta$ 以量词推理结束的情况。
  \begin{enumerate}
    \item 如果最后一个推理是 $\Intro{\lforall}$，则 $\delta$ 形如
    \[
    \AxiomC{}
    \RightLabel{$\delta'$}
    \DeduceC{$!A(a,c)$}
    \RightLabel{\Intro{\lforall}}
    \UnaryInfC{$\lforall[x][!A(x,c)]$}
    \DisplayProof
    \]
    根据归纳假设，$\Subst{\delta'}{t}{c}$ 是 $!A(a,t)$ 的一个正则推导。由于 $a$ 是~$\delta$ 的一个本征变元，$a$~不在~$t$ 中出现。因此，$\Subst{\delta}{t}{c}$ 中的最后一个推理，
    \[
    \AxiomC{}
    \RightLabel{$\Subst{\delta'}{t}{c}$}
    \DeduceC{$!A(a,t)$}
    \RightLabel{\Intro{\lforall}}
    \UnaryInfC{$\lforall[x][!A(x,t)]$}
    \DisplayProof
    \]
    是合法的：因为 $t$ 不包含~$a$，所以结论与任何开放假设均不包含~$a$。
    \item 如果最后一个推理是 $\Elim{\lforall}$，则 $\delta$ 形如
    \[
    \AxiomC{}
    \RightLabel{$\delta'$}
    \DeduceC{$\lforall[x][!A(x,c)]$}
    \RightLabel{\Elim{\lforall}}
    \UnaryInfC{$!A(s(c),c)$}
    \DisplayProof
    \]
    根据归纳假设，$\Subst{\delta'}{t}{c}$ 是 $\lforall[x][!A(x,t)]$ 的一个正则推导。（结论中的项~$s(c)$ 可能包含~$c$，但未必包含。）$\Subst{\delta}{t}{c}$ 中的最后一个推理，
    \[
    \AxiomC{}
    \RightLabel{$\Subst{\delta'}{t}{c}$}
    \DeduceC{$\lforall[x][!A(x,t)]$}
    \RightLabel{\Elim{\lforall}}
    \UnaryInfC{$!A(s(t),t)$}
    \DisplayProof
    \]
    是合法的，且 $!A(s(t),t) \ident \Subst{!A(s(c),c)}{t}{c}$。 \Intro{\lexists} 的情况类似。
  \item 如果最后一个推理是 $\Elim{\lexists}$，则 $\delta$ 形如
    \[
    \AxiomC{}
    \RightLabel{$\delta_1'$}
    \DeduceC{$\lexists[x][!A(x,c)]$}
    \AxiomC{$\Discharge{!A(a,c)}{x}$}
    \RightLabel{$\delta_2'$}
    \DeduceC{$!C$}
    \DischargeRule{\Elim{\lexists}}{x}
    \BinaryInfC{$!C$}
    \DisplayProof
    \]
    根据归纳假设，$\Subst{\delta_1'}{t}{c}$ 和
    $\Subst{\delta_2'}{t}{c}$ 分别是 $\lexists[x][!A(x,c)]$ 的正则推导与从开放假设~$!A(a,t)$ 到 $!C$ 的正则推导。由于 $a$ 是~$\delta$ 的一个本征变元，$a$~不在~$t$ 中出现。因此，$\Subst{\delta}{t}{c}$ 中的最后一个推理，
    \[
    \AxiomC{}
    \RightLabel{$\Subst{\delta'}{t}{c}$}
    \DeduceC{$!A(a,t)$}
    \RightLabel{\Intro{\lforall}}
    \UnaryInfC{$\lforall[x][!A(x,t)]$}
    \DisplayProof
    \]
    是合法的：$\Subst{!A(x,t)}{a}{x} \ident \Subst{!A(a,c)}{t}{c}$，结论~$!C$ 与所有开放假设均不包含~$a$。
  \end{enumerate}
\end{proof}

\begin{prop}\ollabel{prop:regular}
如果 $\delta$ 是 \Log{N1c} 或 \Log{N1i} 中的一个推导，那么存在一个正则推导~$\delta'$，它具有相同的结构（特别是相同的高度），且与 $\delta$ 的区别仅在于替换了本征变元。
\end{prop}

\begin{proof}
仅就本证明而言，如果~$\delta$ 中的一个本征变元推理的本征变元不是其他推理的本征变元，并且除了在该推理的直接子推导中之外，不在~$\delta$ 的任何其他地方出现，则称该推理是\emph{干净的}；否则称其为\emph{脏的}。
令 $n$ 为~$\delta$ 中脏的本征变元推理的个数。
我们将通过对~$n$ 施归纳来证明，$\delta$ 可以转化为所要求的正则推导~$\delta'$。
如果 $n=0$，则 $\delta$ 中的所有本征变元推理都是干净的，因此 $\delta$ 是正则的，无需再证。
否则，选取一个最顶端的本征变元推理，设为一个 $\Intro{\lforall}$ 推理。
$\delta$ 中以它结尾的子推导~$\delta_1$ 形如
    \[
    \AxiomC{}
    \RightLabel{$\delta_2$}
    \DeduceC{$!A(c)$}
    \RightLabel{\Intro{\lforall}}
    \UnaryInfC{$\lforall[x][!A(x)]$}
    \DisplayProof
    \]
注意，由于 $c$ 是本征变元，它不出现在~$\lforall[x][!A(x)]$ 或任何开放假设中。
推导 $\delta_2$ 是正则的，因为它根本不包含本征变元推理。
令 $a$ 为一个不在~$\delta$ 中出现的常项符号。
根据 \olref{lem:subst-regular}，$\Subst{\delta_2}{a}{c}$ 是 $!A(a)$ 的一个正则推导。
根据本征变元条件，$\delta_2$ 的开放假设不包含~$c$，
因此~$\Subst{\delta_2}{a}{c}$ 的开放假设也不包含~$a$。
于是我们可以应用~$\Intro{\lforall}$。
如果将 $\delta$ 中的 $\delta_1$ 替换为推导 $\delta_1'$：
    \[
    \AxiomC{}
    \RightLabel{$\Subst{\delta_2}{a}{c}$}
    \DeduceC{$!A(a)$}
    \RightLabel{\Intro{\lforall}}
    \UnaryInfC{$\lforall[x][!A(x)]$}
    \DisplayProof
    \]
我们就将一个脏的本征变元推理替换成了一个干净的，唯一的区别是将本征变元 $c$ 换成了~$a$。
所得的推导含有 $n-1$ 个脏的本征变元推理，
从而由归纳假设即得结论。
\end{proof}

\begin{prob}
描述 \olref{prop:regular} 的证明中，最顶端的本征变元推理是~\Elim{\lexists} 的情形。
\end{prob}

我们不会显式地引用刚刚证明的这个命题，但我们将假设所考虑的所有推导都是正则的——如果它们不是，
我们总可以通过替换本征变元使其成为正则的。

\Olref{lem:subst-regular} 和 \olref{prop:regular} 对于 \Log{N2c} 和~\Log{N2i} 也成立。
我们将证明留作练习。

\end{document}